%% Versión resumida (una carilla) del curriculum en español.
%% Basado en curriculum-español.tex.

\documentclass[10pt,a4paper,sans]{moderncv}

\moderncvstyle{casual}
\moderncvcolor{blue}
\nopagenumbers{}

\usepackage[utf8]{inputenc}

% márgenes chicos para que entre en una carilla
\usepackage[top=1.2cm,bottom=1.2cm,left=1.5cm,right=1.5cm]{geometry}
\setlength{\hintscolumnwidth}{3cm}

% títulos de sección más chicos y con menos espacio
\renewcommand*{\sectionfont}{\large\mdseries\upshape}
\makeatletter
\settoheight{\baseletterheight}{\sectionstyle{o}}
\setlength{\baseletterheight}{\baseletterheight-0.95ex}
\RenewDocumentCommand{\section}{sm}{%
  \par\addvspace{1ex}%
  \phantomsection{}%
  \addcontentsline{toc}{section}{#2}%
  \parbox[t]{\hintscolumnwidth}{\strut\raggedleft\raisebox{\baseletterheight}{\color{color1}\rule{\hintscolumnwidth}{0.95ex}}}%
  \hspace{\separatorcolumnwidth}%
  \parbox[t]{\maincolumnwidth}{\strut\sectionstyle{#2}}%
  \par\nobreak\addvspace{0.4ex}\@afterheading}
\makeatother

% entrada compacta de una línea: \cvcompact{fechas}{cargo}{lugar}
\newcommand{\cvcompact}[3]{\cvitem[0.2ex]{\footnotesize #1}{\small\textbf{#2}, #3}}
% ítem de lista compacto
\newcommand{\cvsmallitem}[1]{\cvlistitem[0.2ex]{\small #1}}

\name{Eric}{Brandwein}
\phone[mobile]{+54~(911)~6120~4615}
\email{ebrandwein@dc.uba.ar}
\homepage{ericbrandwein.github.io}

\begin{document}
\makecvtitle
\vspace{-4ex}

\section{Estudios}
\cvcompact{2023 -- 2027 (est.)}{Doctorado en Ciencias de la Computación}{Universidad de Buenos Aires. Directores: Flavia Bonomo e Ignasi Sau. Tema: Kernelización en grafos.}
\cvcompact{2015 -- 2022}{Licenciatura en Ciencias de la Computación}{Universidad de Buenos Aires.}

\section{Docencia}
\cvcompact{03/2026 -- 06/2026}{Profesor Titular}{Matemática Discreta, Ing. en Inteligencia Artificial, Universidad del CEMA.}
\cvcompact{10/2023 -- 08/2026}{Ayudante de 1ra / JTP / Cargo Equiparado a JTP}{Algoritmos y Estructuras de Datos III, DC, FCEyN, UBA.}
\cvcompact{08/2024 -- 02/2025}{Profesor Asistente}{Algoritmos y Estructura de Datos, Depto.\ de Ingeniería, Universidad de San Andrés.}
\cvcompact{02/2019 -- 02/2020}{Ayudante de 2da}{Algoritmos y Estructuras de Datos I, DC, FCEyN, UBA.}

\section{Publicaciones}
\cvsmallitem{\emph{Kernelization dichotomies for hitting minors under structural parameterizations}. M.~Bougeret, E.~Brandwein, I.~Sau. STACS 2026, LIPIcs, 17:1--17:19.}
\cvsmallitem{\emph{On the thinness of trees}. F.~Bonomo-Braberman, E.~Brandwein, C.~L.~Gonzalez, A.~Sansone. Discrete Applied Mathematics 365:39--60, 2025. Versión preliminar en ISCO 2022, LNCS 13526.}
\cvsmallitem{\emph{Thinness 2 condicionada a órdenes de co-bipartición}. A.~Dinkel, F.~Bonomo-Braberman, E.~Brandwein. JAIIO 11(14):98--111, 2025.}
\cvsmallitem{\emph{Thinness and its variations on some graph families and coloring graphs of bounded thinness}. F.~Bonomo-Braberman, E.~Brandwein, F.~de~S.~Oliveira, M.~S.~Sampaio~Jr., A.~Sansone, J.~L.~Szwarcfiter. RAIRO-Operations Research 58(2):1681--1702, 2024.}
\cvsmallitem{\emph{Computing parameters that generalize interval graphs using restricted modular partitions}. F.~Bonomo-Braberman, E.~Brandwein, I.~Sau. Enviado para revisión.}

\section{Charlas}
\cvsmallitem{\emph{Kernelización}. 8º Día de la Investigación en Ciencias de la Computación, ICC, UBA, Ago. 2026.}
\cvsmallitem{\emph{Dicotomías de kernelización para minor hitting}. I Encuentro Invernal de Grafos, UNSL, San Luis, Jul. 2026.}
\cvsmallitem{\emph{Kernelization dichotomies for hitting minors under structural parameterizations}. Séminaire ALGCo, LIRMM, Montpellier, Francia, Oct. 2025.}
\cvsmallitem{\emph{Parameterized algorithms for thinness via the cluster module number}. LAWCG 2024, Ceará, Brasil, Oct. 2024; Workshop SINFIN, Buenos Aires, Mar. 2024.}

\section{Tesistas}
\cvitem[0.2ex]{\footnotesize Lic.\ en Computación}{\small Mauricio Romero Laino, Ignacio Lucas Rodriguez, Ezequiel Maidanik, Bárbara Rivero, Manuela Martínez (co-dir.).}
\cvitem[0.2ex]{\footnotesize Lic.\ en Datos}{\small Germán Rud, Ayelén Dinkel (co-dir., defendida Dic.\ 2024).}
\cvitem[0.2ex]{\footnotesize Lic.\ en Matemática}{\small Manuel Robert (co-dir., defendida Nov.\ 2025).}

\section{Becas}
\cvsmallitem{\textbf{Beca de doctorado} CONICET (2024 -- 2029) y UBACyT (2023 -- 2024).}
\cvsmallitem{\textbf{Visitas de investigación} al LIRMM, Montpellier, Francia (Ago. 2024 y Oct. 2025). \textbf{+Viajes 2024}, UBA.}

\section{Otras actividades}
\cvitem[0.2ex]{\footnotesize Revisión}{\small Revisor para LAGOS 2025, ISAAC 2024 e IPEC 2024.}
\cvitem[0.2ex]{\footnotesize Comités}{\small Comité organizador de LAGOS 2025.}

\section{Experiencia en la industria}
\cvitem[0.2ex]{\footnotesize 2015 -- 2022}{\small Ingeniero de Software en Wildlife Studios (2022), Mulesoft (2021), Medallia (2021) y Mercadolibre (2015--2017).}

\section{Idiomas y otros}
\cvitem[0.2ex]{\footnotesize Idiomas}{\small Español (nativo), inglés (fluido, CAE Cambridge), italiano (intermedio), francés (básico).}
\cvitem[0.2ex]{\footnotesize Competencias}{\small Tercer puesto en el TAP 2019 y 2020. 11 medallas y reconocimientos en OIA y OMA, incluyendo medalla de oro nacional.}

\end{document}
